\begingroup
\setlength{\parskip}{4pt}
\setlength{\tabcolsep}{3pt}
\renewcommand{\arraystretch}{1.40}
\setstretch{1}
\footnotesize
\noindent\textbf{FICHA 02}\par
{\bfseries Ficha de estructuras de drenaje transversal}\par
Corredor EMP.LI105 Pedregal--Censo--Cumbray, Simbal | Evaluación 2026\par
Mida el área interior y la superficie ocupada en la entrada y la salida del pase.\par
\begin{tabularx}{\linewidth}{|p{2.1cm}|p{1.1cm}|X|p{3.3cm}|}
\hline
\textbf{Código unidad} & \textbf{Tramo} & \textbf{Progresivas inicio y fin} & \textbf{Fecha y hora de medición}\\\hline
\rule{0pt}{4mm} & & & \\\hline
\end{tabularx}
Evaluador \rule{3.3cm}{0.25pt}\quad Pase de agua \rule{2.9cm}{0.25pt}\par
Coordenadas WGS84 \rule{2.4cm}{0.25pt}\quad Lluvia o mantenimiento \rule{2.2cm}{0.25pt}\par
Registre las dimensiones en metros. Para una abertura irregular, use una fotografía frontal con escala y delimite el contorno mediante un polígono. Registre área libre igual a cero cuando toda la abertura esté ocupada.\par
\begin{tabularx}{\linewidth}{|p{1.05cm}|X|p{1.35cm}|p{1.35cm}|p{1.25cm}|p{1.35cm}|p{0.75cm}|}
\hline
\textbf{Extremo} & \textbf{Forma y dimensiones (m)} & \textbf{Área interior (m$^2$)} & \textbf{Área ocupada (m$^2$)} & \textbf{Área libre (m$^2$)} & \textbf{Abertura libre (\%)} & \textbf{Foto}\\\hline
Entrada&&&&&&\\[2pt]\hline
Salida&&&&&&\\[2pt]\hline
\end{tabularx}
Área libre = área interior $-$ área ocupada. Abertura libre (\%) = 100 $\times$ área libre / área interior.\par
CT (\%) = menor porcentaje de abertura libre entre entrada y salida. Resultado \rule{1.8cm}{0.25pt} \%\par
\begin{tabularx}{\linewidth}{|p{4cm}|X|}
\hline
\textbf{Aspecto complementario} & \textbf{Descripción y evidencia}\\\hline
Integridad de la estructura & \rule{0pt}{5mm}\\\hline
Erosión & \rule{0pt}{5mm}\\\hline
Condición de descarga & \rule{0pt}{5mm}\\\hline
\end{tabularx}
Instrumentos auxiliares: regla graduada, wincha y cámara.\par
Observaciones \rule{9.0cm}{0.25pt}\par
Responsable \rule{3.8cm}{0.25pt}\qquad Firma \rule{3.1cm}{0.25pt}
\endgroup
